%%%%%%%%%%%%%%%%%%%%%%%%
%
% $Autor: Müller, Hinrichs $
% $Datum: 2026-06-09 14:18:35Z $
% $Pfad: UltraschallradarAutomatisierungstechnik/report/System/Ultraschallradar/Radar.tex $
% $Version: 1 $
%
%
%%%%%%%%%%%%%%%%%%%%%%%%

\part{Einleitung} 

\chapter{Einführung}

Die Umgebungserfassung und die präzise Detektion von Hindernissen sind fundamentale Bestandteile moderner Automatisierungs-, Automobil- und Robotiksysteme. Ein klassisches, aber nach wie vor anschauliches Anwendungsbeispiel hierfür ist das Radar-Prinzip beziehungsweise das hier adaptierte Prinzip des Sonars (Sound Navigation and Ranging). 

Im Rahmen dieses Projekts wird ein voll funktionsfähiges Ultraschallradar konzipiert, mechanisch konstruiert und informationstechnisch implementiert. Ziel der Arbeit ist es, das komplexe Zusammenspiel von Mikrocontroller-Hardware, Sensorik, Aktuatorik und einer grafischen Benutzeroberfläche an einem praxisnahen, mechatronischen System zu demonstrieren.

Als zentrale Steuereinheit des Systems dient ein Arduino Nano. Die mechanische Scan-Bewegung wird durch einen 180\textdegree{}-Servomotor vom Typ SG90 realisiert, auf dessen Rotationsachse ein HC-SR-04 Ultraschallsensor montiert ist. Um eine stabile Plattform für diese Komponenten zu gewährleisten, wurde im Zuge des Projekts ein ma\ss passendes Gehäuse im 3D-Druck-Verfahren gefertigt. 

Auf der Software-Seite gliedert sich das Projekt in zwei Hauptbereiche: Die Programmierung des Mikrocontrollers (Embedded C/C++) über die Arduino IDE sowie die Entwicklung einer Benutzeroberfläche für den Computer. Letzteres wird mithilfe der auf Java basierenden Umgebung Processing umgesetzt. Diese Software empfängt die seriellen Sensordaten des Arduinos (Winkel und Distanz) kontinuierlich und übersetzt sie in eine dynamische, visuell ansprechende Radar-Anzeige, welche erkannte Objekte und Systemzustände visualisiert.

Der vorliegende Bericht dokumentiert den gesamten Entwicklungsprozess dieses Systems. Er erläutert zunächst die verwendeten Hardwarekomponenten, die mechanische Konstruktion sowie den elektronischen Schaltungsaufbau. Anschließend wird detailliert auf die Programmierung und Implementierung der Steuerungs- sowie Visualisierungssoftware eingegangen.

\chapter{Aufgabenstellung}

	Im Rahmen der Veranstaltung Automatisierungstechnik soll eine Applikation entwickelt werden, die die Funktionsfähigkeit des Arduino Nano 33 BLE Sense Lite sowie ausgewählter Sensoren und Aktoren demonstriert. Die vorliegende Arbeit befasst sich mit der Realisierung eines Systems zur Umgebungserfassung mittels eines Ultraschall-Abstandssensors, der auf einem Servomotor montiert ist und von einem Arduino Nano 33 BLE Sense Lite angesteuert wird. Ziel ist die Umsetzung eines radarartigen Scans der Umgebung und die Visualisierung der gemessenen Distanzen in Form einer geeigneten Darstellung auf dem PC.
	
	Konkret umfasst die Aufgabenstellung die Integration eines Ultraschall-Abstandssensors (HC-SR-04) und eines Servomotors (Tower Pro SG90) in die vorhandene Hardwareplattform. Der Servomotor soll den Sensor in einem definierten Winkelbereich schwenken, während der Mikrocontroller in vorgegebenen Winkelinkrementen Entfernungswerte aufnimmt und diesen Messdaten einen eindeutigen Winkel zuordnet. Auf Basis dieser Messwerte ist eine Auswertung zu konzipieren, die eine übersichtliche Darstellung der Umgebung ermöglicht.
	
	Neben der reinen Funktionserfüllung werden weitere Anforderungen an die
	Applikation gestellt. Dazu gehören eine strukturierte und modular aufgebaute Firmware, die Verwendung geeigneter Bibliotheken zur Ansteuerung von Sensor und Aktor. Au\ss erdem sind Verfahren zur Visualisierung der Live-Daten auf einem PC zu entwickeln und zu implementieren, sodass die resultierende Umgebungskarte in Echtzeit nachvollzogen werden kann.
	
	Die Aufgabenstellung umfasst weiterhin die Erstellung einer Entwickler- und Service-Dokumentation. Diese beinhaltet neben der Projektbeschreibung eine detaillierte Beschreibung der Hardware des Arduino Nano 33 BLE Sense Lite und der eingesetzten Komponenten, die Spezifikation und Durchführung von Tests, eine Darstellung der Softwarearchitektur sowie Angaben zu Bezugsquellen und Stückliste (Bill of Materials).
	
\chapter{Stand der Technik}

	Ultraschall-Abstandssensoren werden seit vielen Jahren in der Automatisierungstechnik, Robotik und Fahrzeugtechnik zur berührungslosen Distanzmessung eingesetzt.\cite{leuze:2026} Sie arbeiten in der Regel nach dem Laufzeitmessprinzip, bei dem ein kurzer Ultraschallimpuls ausgesendet und die Zeit bis zum Empfang des reflektierten Echos gemessen wird. Aus der Schalllaufzeit lässt sich unter Annahme einer bekannten Schallgeschwindigkeit der Abstand zu einem Hindernis bestimmen. Im Niedrigkostenbereich haben sich dafür Sensormodule wie der HC-SR-04 etabliert, die eine einfache Anbindung an Mikrocontroller über Trigger- und Echo-Signale erlauben.
	
	Für die mechanische Ausrichtung von Sensoren und leichten Lasten werden in zahlreichen Anwendungen kompakte RC‑ Servomotoren eingesetzt. Diese enthalten neben einem Gleichstrommotor ein Getriebe sowie eine integrierte Positionsregelung, die über ein Pulsweitenmodulationssignal (PWM) angesteuert wird. Miniaturservos wie der Tower Pro SG90 sind vor allem in ferngesteuerten Modellen wie Flugzeugen und Helikoptern verbreitet, wo sie etwa Ruderflächen oder andere bewegliche Komponenten ansteuern.\cite{kts:2026} In technischen Lehr- und Demonstrationsaufbauten werden derartige Servos zudem häufig genutzt, um Sensoren in einem begrenzten Winkelbereich zu positionieren.
	
	Die Kombination aus Ultraschall-Abstandssensor und Servomotor ermöglicht die Realisierung von Systemen zur zweidimensionalen Umgebungserfassung. Durch das sequenzielle Abtasten verschiedener Winkelpositionen kann ein Scan der Umgebung erzeugt werden. In frei verfügbaren Projekten und Tutorials wird dieses Prinzip typischerweise mithilfe eines Ultraschallmoduls des Typs HC-SR-04, eines SG90-Servos und eines Arduino Uno oder Arduino Nano umgesetzt. (vgl. \cite{projehocam:2015}) Die gemessenen Distanzen werden dabei entweder auf einem seriell angebundenen PC visualisiert oder direkt auf einem Display dargestellt. Derartige Beispiele dienen häufig als Einführung in die Sensordatenerfassung und -visualisierung mit einfachen Mikrocontrollersystemen.
	


\chapter{Einbettung des Projektes}


	Das in diesem Projekt realisierte System knüpft an den beschriebenen Stand der Technik an, indem es die etablierte Kombination aus Ultraschall-Abstandssensor HC-04 und Tower Pro SG90-Servomotor zur Umgebungserfassung verwendet. Der Ultraschallsensor ist auf dem Servomotor montiert und wird in einem definierten Winkelbereich geschwenkt, sodass in diskreten Winkelpositionen Entfernungswerte aufgenommen und zu einem Umgebungsbild zusammengeführt werden.
	
	Viele veröffentlichte Beispielprojekte setzen noch auf Arduino-Boards mit 8‑Bit‑Mikrocontrollern.\cite{Kaparthi:2021} In der vorliegenden Applikation kommt dagegen der Arduino Nano 33 BLE Sense als Steuerplattform zum Einsatz. Diese Hardwarebasis ermöglicht eine strukturierte Implementierung der Mess- und Auswertealgorithmen sowie die Erweiterung des klassischen Ultraschall-Radar-Konzepts um weitergehende Funktionen, etwa eine flexible Visualisierung der Messdaten auf einem PC.
	
	Durch diese Ausgestaltung ordnet sich das Projekt in die Gruppe kostengünstiger Ultraschall-Radar-Systeme ein, entwickelt diese jedoch durch den Einsatz moderner Hardware und eine systematische Softwarestruktur weiter.\cite{Kaparthi:2021} Es dient damit sowohl als Demonstrator für die Sensordatenerfassung mit dem Arduino Nano 33 BLE Sense als auch als Grundlage für weiterführende Anwendungen in der Automatisierungstechnik.

\chapter{Besondere Herausforderung}

	Im Rahmen der Entwicklung der Applikation traten mehrere besondere Herausforderungen auf, die sowohl die mechanische Konstruktion als auch die Softwaregestaltung betrafen. Ein zentrales Ziel war der möglichst kompakte Aufbau des Systems. Dabei musste das Gehäuse so dimensioniert und gestaltet werden, dass alle Komponenten, insbesondere der Ultraschallsensor, der Servomotor und das Arduino‑Board, platzsparend integriert werden konnten, ohne die Funktionalität oder Zugänglichkeit zu beeinträchtigen. Eine zusätzliche Herausforderung bestand in der Unterbringung der Jumper-Kabel im Inneren des Gehäuses, um mechanische Kollisionen, Kabelbrüche sowie Beeinflussungen der Sensorik durch bewegte Leitungen zu vermeiden.
	
	Auf der Softwareseite ergaben sich weitere Anforderungen aus der gewünschten Plattformunabhängigkeit der Benutzerschnittstelle. Es sollte eine Anwendung bereitgestellt werden, die sich direkt von einem USB ‑ Stick starten lässt und möglichst ohne zusätzliche Installation auf unterschiedlichen Betriebssystemen lauffähig ist. Dies erforderte eine genaue Auswahl der eingesetzten Entwicklungsumgebung, um Abhängigkeiten zu reduzieren und eine robuste Ausführung zu ermöglichen.
	
	Im Verlauf der Entwicklung zeigte sich zudem, dass die Visualisierung der Messdaten auf der Benutzeroberfläche zunächst zu grob und unübersichtlich war. Einzelne Objekte wurden zu gro\ss dargestellt, sodass die eigentlichen Objektdimensionen anhand des Ultraschallradars nur eingeschränkt erkennbar waren.
	
	
\section{Lösung der Herausforderung}


\chapter{Inhalte der folgenden Kapitel}

